\section{Independent original-source control of both discarded intervals}

This proof uses the original relation polynomial, the original Gamma density,
and the original polynomial coefficient spaces.  The change of variable below
is an explicit evaluation map, with its complete coefficient matrix.  No
mass factor or polynomial degree is removed from a Gram matrix.

\subsection{Objects, intervals, and the coefficient map}

Fix integers $m\ge1$ and $l\ge1$, and put
\begin{equation}
 k=4l+1,\quad c=k/2,\quad e=1+k(m-1),\quad q=e(k+1)^2,
 \quad 0<\delta<\tfrac12,\quad \gamma>2.
 \tag{FTC1}
\end{equation}
Write $d=\delta^2$ and $g_\gamma=\gamma^2$.  The notation $g_\gamma$
does not reassign the arithmetic entire function denoted by $g$ elsewhere.
Define
\begin{align}
 w_{ab}&=(2a-k)\delta+i(2b-k)\gamma,\qquad 0\le a,b\le k,\notag\\
 \chi(S)&=\prod_{a,b=0}^k(S-c-w_{ab})^e,
 &R_{\max}&=k\sqrt{d+g_\gamma},\notag\\
 t_j&=2j+\tfrac12,\quad
 f_l(S)=\prod_{j=0}^{l-1}(S-c-t_j),
 &t_{\max}&=2l-\tfrac32,\notag\\
 \sigma(y)&=\frac{|\Gamma(\tfrac14+iy/2)|^2}{2\pi},
 &\phi_0(S)&=1,\qquad \phi_l(S)=f_l(S).
 \tag{FTC2}
\end{align}
Thus $|w_{ab}|\le R_{\max}$, the multiset of roots $w_{ab}$ is invariant
under $w\mapsto-w$, and
$|\phi_l(c+iy)|^2=\prod_{j=0}^{l-1}(y^2+t_j^2)$.
The two relation dimensions treated here are $r=q$ and $r=q+1$.
For $P(S)=\sum_{n=0}^{r-1}p_nS^n$ and a Borel set $A\subset\mathbb R$
define the original form
\begin{equation}
 \mathcal H_{s,A}(P,P)=
 \int_{\{y:y/q\in A\}}
 |P(c+iy)\chi(c+iy)\phi_s(c+iy)|^2\sigma(y)\,dy,
 \qquad s\in\{0,l\}.
 \tag{FTC3}
\end{equation}
Its matrix $H_{s,A}$ is in the unchanged ordered basis
$(1,S,\ldots,S^{r-1})$; in particular its $(i,j)$ entry, indexed from zero,
contains $\overline{(c+iy)^i}(c+iy)^j$.  Set
\begin{equation}
 Q(z)=P(c+iqz),\qquad
 (T_{c,q})_{jn}=\begin{cases}
 \binom nj c^{n-j}(iq)^j,&j\le n,\\0,&j>n.
 \end{cases}
 \quad \det T_{c,q}=(iq)^{r(r-1)/2}.
 \tag{FTC4}
\end{equation}
The coefficient vector of $Q$ is exactly $T_{c,q}p$.  If $\widetilde H_{s,A}$
is the matrix in powers of $z$ for the measure obtained by the literal
substitution $y=qz$, including the Jacobian $q$, then
$H_{s,A}=T_{c,q}^*\widetilde H_{s,A}T_{c,q}$ and
$\det H_{s,A}=q^{r(r-1)}\det\widetilde H_{s,A}$.
All subsequent comparisons are comparisons of the original matrices
$H_{s,A}$.  These congruences prove their exact relation to the intermediate
$z$ coordinate calculation.

\subsection{Gamma bounds with the original mass}

For $0<\theta<\pi/2$ and every real $y$,
\begin{equation}
 \frac{\pi}{\Gamma(3/4)^2}e^{-\pi|y|}
 \le\sigma(y)\le
 \frac{\Gamma(1/4)^2}{2\pi\sqrt{\cos\theta}}
 e^{-\theta|y|}.
 \tag{FTC5}
\end{equation}
Here is a proof of the constants.  For $a>0$ and $b\ge0$, rotate the
Euler integral for $\Gamma(a+ib)$ through the angle $\theta$ in the first
quadrant, using the branch of the logarithm with argument in $[0,\theta]$.
The small circular arc is $O(\varepsilon^a)$, and the large arc tends to
zero because $\cos v\ge\cos\theta>0$ on the rotation sector.  Consequently
\[
 \Gamma(a+ib)=e^{i\theta(a+ib)}
 \int_0^\infty t^{a+ib-1}e^{-e^{i\theta}t}\,dt,
 \qquad
 |\Gamma(a+ib)|\le e^{-\theta b}
 (\cos\theta)^{-a}\Gamma(a).
\]
For $b<0$ use the conjugate sector.  Squaring with $a=1/4$, $b=y/2$
proves the upper bound in (FTC5).

For completeness, the reflection identity used for the lower bound follows
from Euler's double integral, with the substitution $u=vt$ followed by
integration in $v$:
\[
 \Gamma(z)\Gamma(1-z)=\int_0^\infty\frac{t^{z-1}}{1+t}\,dt,
 \qquad 0<\operatorname{Re}z<1.
\]
On a keyhole contour with $0<\arg t<2\pi$, the integrals over the small
and large circles vanish because $0<\operatorname{Re}z<1$.  The two
straight integrals have difference $(1-e^{2\pi iz})$ times the displayed
integral, whereas the residue at $-1$ is $-e^{\pi iz}$.  The residue formula
therefore gives $\Gamma(z)\Gamma(1-z)=\pi/\sin(\pi z)$.
The unrotated Euler integral gives
$|\Gamma(3/4-iy/2)|\le\Gamma(3/4)$, and
\[
 |\sin(\pi/4+i\pi y/2)|^2=\tfrac12\cosh(\pi y).
\]
It follows that
\[
 \sigma(y)\ge\frac{\pi}{\Gamma(3/4)^2\cosh(\pi y)}
 \ge\frac{\pi}{\Gamma(3/4)^2}e^{-\pi|y|}.
\]
In particular the exact ratio of the upper and lower constant in (FTC5)
is $(\cos\theta)^{-1/2}$, since the same reflection calculation at $z=1/4$
gives $\Gamma(1/4)\Gamma(3/4)=\pi\sqrt2$.  Denote the lower constant by
$c_L=\pi/\Gamma(3/4)^2$.

\subsection{Exterior expansion of the unchanged relation polynomial}

Define
\begin{equation}
 a_k=\frac{k(k+2)(d-g_\gamma)}{3q}<0,\qquad
 \tau_x=\frac{k^2(d+g_\gamma)}{q^2x^2},\qquad
 \Delta_\chi(x)=\frac{q\tau_x^2}{2(1-\tau_x)}
 \quad(\tau_x<1).
 \tag{FTC6}
\end{equation}
For real $|z|\ge x$ with $\tau_x<1$ the exact logarithmic expansion is
\begin{equation}
 \log\frac{|\chi(c+iqz)|^2}{(q|z|)^{2q}}
 =\frac{a_k}{z^2}+E_k(z),\qquad |E_k(z)|\le\Delta_\chi(x).
 \tag{FTC7}
\end{equation}
Indeed the series for $\log(1-w_{ab}/(iqz))$ converges absolutely and
uniformly on this real exterior set.  Pairing $w$ and $-w$ cancels every
odd power.  Its degree-two term is
\[
 \frac{e}{q^2z^2}\sum_{a,b=0}^k\operatorname{Re}(w_{ab}^2)
 =\frac{e(k+1)^2k(k+2)(d-g_\gamma)}{3q^2z^2}
 =\frac{a_k}{z^2},
\]
where $\sum_{a=0}^k(2a-k)^2=k(k+1)(k+2)/3$ follows by expanding the
square and summing $a$ and $a^2$.  The remaining series is
\[
 E_k(z)=-e\operatorname{Re}
 \sum_{a,b=0}^k\sum_{n=2}^\infty
 \frac1n\left(\frac{w_{ab}}{iqz}\right)^{2n}.
\]
Thus its absolute value is at most
$q\sum_{n=2}^\infty\tau_x^n/n\le q\tau_x^2/[2(1-\tau_x)]$.
Also put
\begin{equation}
 \Delta_{f,0}(x)=0,\qquad
 \Delta_{f,l}(x)=\frac{\sum_{j=0}^{l-1}t_j^2}{q^2x^2}.
 \quad
 (q|z|)^{2s}\le|\phi_s(c+iqz)|^2
 \le(q|z|)^{2s}e^{\Delta_{f,s}(x)}\quad(|z|\ge x).
 \tag{FTC8}
\end{equation}
The last upper estimate is exactly $\log(1+u)\le u$ applied separately
to every factor; none of the $t_j$ have been changed.

\subsection{A polynomial estimate on the original coefficient space}

For every complex polynomial $Q$ of degree at most $r-1$, write
$I_Q=\int_1^2|Q(t)|^2dt$.  For real $|z|\ge B\ge2$,
\begin{equation}
 |Q(z)|^2\le r^2\bigl((4+6/B)|z|\bigr)^{2(r-1)}I_Q.
 \tag{FTC9}
\end{equation}
For $0<\varepsilon\le1/2$ and real $|z|\le\varepsilon$,
\begin{equation}
 |Q(z)|^2\le r^2(6+4\varepsilon)^{2(r-1)}I_Q.
 \tag{FTC10}
\end{equation}
To prove both statements, use the polynomials
$\psi_n(t)=\sqrt{2n+1}\,P_n(2t-3)$, where
\[
 P_n(x)=\frac{1}{2^n n!}\frac{d^n}{dx^n}(x^2-1)^n.
\]
Integration by parts proves orthogonality to every lower-degree polynomial:
all boundary derivatives of $(x^2-1)^n$ of order less than $n$ vanish.
The leading coefficient is $(2n)!/[2^n(n!)^2]$, and the same integration
by parts, together with
\[
 \int_{-1}^1(1-x^2)^n dx
 =\frac{2^{2n+1}(n!)^2}{(2n+1)!},
\]
gives $\int_{-1}^1P_n(x)^2dx=2/(2n+1)$.  The last integral follows from
$x=2t-1$ and $n$ integrations by parts for
$\int_0^1t^n(1-t)^n dt=(n!)^2/(2n+1)!$.
Hence $(\psi_n)_{0\le n<r}$ is an orthonormal basis on $[1,2]$.

For real $x\ge1$ the finite integral formula
\[
 P_n(x)=\frac1\pi\int_0^\pi
 \bigl(x+\sqrt{x^2-1}\cos v\bigr)^n dv
\]
is obtained by expanding the integrand: odd cosine moments vanish and
the $2j$-th moment is $\binom{2j}{j}/4^j$, giving the same polynomial as
the displayed Rodrigues derivative.  It implies $|P_n(x)|\le(2x)^n$.
Parity of Rodrigues' derivative gives this bound with $2|x|$ for
$x\le-1$ as well.  Expansion of $Q$ in the orthonormal basis and complex
Cauchy--Schwarz now give
\[
 |Q(z)|^2\le I_Q\sum_{n=0}^{r-1}(2n+1)|P_n(2z-3)|^2.
\]
On the exterior set, $|2z-3|\ge1$ and
$2|2z-3|\le4|z|+6\le(4+6/B)|z|$.  On the inner interval,
$|2z-3|\ge2$ and $2|2z-3|\le6+4\varepsilon$.
Each bound on the right is at least one, and
$\sum_{n=0}^{r-1}(2n+1)=r^2$.  This proves (FTC9)--(FTC10), including
their degree and constants.

\subsection{Far interval: a finite bound and a uniform exponential bound}

Let $\mathcal B_B=\{z:1\le|z|\le B\}$ and
$\mathcal F_B=\{z:|z|>B\}$, where $B\ge2$.
Suppose first that $\tau_1<1$, and set
$\eta=-a_k+\Delta_\chi(1)$.  The positive segment $[1,2]$ gives
\begin{equation}
 \mathcal H_{s,\mathcal B_B}(P,P)
 \ge q^{2q+2s+1}c_L e^{-2\pi q-\eta} I_Q.
 \tag{FTC11}
\end{equation}
In fact $a_k/z^2\ge a_k$ on this segment, the factor $z^{2q+2s}$ is at
least one, and (FTC5), (FTC7), (FTC8), and the Jacobian $dy=q\,dz$ give
exactly this inequality.

Write $D_s=2q+2s+2r-2$, so that $D_s\ge0$.  If
$0<\theta<\pi/2$ and $\theta q>D_s/B$, then
\begin{align}
 \mathcal H_{s,\mathcal F_B}&\le\kappa_{s,r}^{\mathrm{far}}(B,\theta)
 \mathcal H_{s,\mathcal B_B},\notag\\
 \kappa_{s,r}^{\mathrm{far}}(B,\theta)
 &=\frac{2r^2}{\sqrt{\cos\theta}}
 e^{2\pi q+\eta+\Delta_\chi(B)+\Delta_{f,s}(B)}
 (4+6/B)^{2(r-1)}
 \frac{B^{D_s}e^{-\theta qB}}{\theta q-D_s/B}.
 \tag{FTC12}
\end{align}
These are quadratic-form inequalities on the full original relation space.
To verify the numerator, use $a_k/z^2\le0$ in the far interval, then
(FTC5), (FTC7)--(FTC9).  The integral over its two components is bounded
by
\[
 \frac{2\Gamma(1/4)^2}{2\pi\sqrt{\cos\theta}}
 q^{2q+2s+1}r^2(4+6/B)^{2(r-1)}
 e^{\Delta_\chi(B)+\Delta_{f,s}(B)}I_Q
 \int_B^\infty z^{D_s}e^{-\theta qz}dz.
\]
For $z\ge B$, $D_s\log(z/B)\le D_s(z-B)/B$ because $D_s\ge0$.
Integration of the resulting exponential proves
\[
 \int_B^\infty z^{D_s}e^{-\theta qz}dz
 \le\frac{B^{D_s}e^{-\theta qB}}{\theta q-D_s/B}.
\]
Divide by (FTC11) and use the exact Gamma-constant ratio proved after
(FTC5).  This proves (FTC12).

For $B=64$, $\theta=1$, the original degrees give
$D_s\le4q+2l\le5q$, since $r-1\le q$ and $2l\le q$.
Therefore
\begin{align}
 \kappa_{s,r}^{\mathrm{far}}(64,1)
 &\le\frac{128r^2}{59q\sqrt{\cos1}}
 e^{\eta+\Delta_\chi(64)+\Delta_{f,s}(64)-c_*q},\notag\\
 c_*&=64-2\pi-5\log64-2\log(131/32)>20.
 \tag{FTC13}
\end{align}
Here $4+6/64=131/32$, $q-D_s/64\ge59q/64$, and $r-1\le q$.
The strict numerical bound uses $\pi<4$, $\log2<1$, and $131/32<8$,
which give $2\pi+5\log64+2\log(131/32)<8+30+6=44$.

There is an explicit finite domain on which this becomes uniform in
the original packet parameters:
\begin{equation}
 q\ge2k\sqrt{d+g_\gamma}
 \quad\Longrightarrow\quad
 \kappa_{s,r}^{\mathrm{far}}(64,1)
 \le\frac{512q}{59\sqrt{\cos1}}e^{-19q}.
 \tag{FTC14}
\end{equation}
Indeed $\tau_1\le1/4$ and $k\ge5$ give, separately,
\[
 -a_k\le\frac{k+2}{3k}q\tau_1\le\frac{7q}{60},\qquad
 \Delta_\chi(1)\le\frac q{24},\qquad
 \Delta_\chi(64)\le\frac q{24}.
\]
Their sum is at most $q/5$.  Since $t_j<2l$,
$\sum_jt_j^2\le4l^3$ and $2l\le q$ imply
$\Delta_{f,l}(64)\le q/8192$.  Thus the sum of all three exponent
corrections in (FTC13) is strictly less than $q/4$.
Equation (FTC13), $c_*>20$, and $r\le q+1\le2q$ prove (FTC14).
The threshold is a stated domain for a finite explicit bound; for each
fixed original $\delta,\gamma,m$ it holds for all sufficiently large
$k=4l+1$, because $q\ge k^2$.

\subsection{Inner interval: the exact original-root product bound}

Set $\varepsilon=2^{-10}$, $\mathcal I_\varepsilon=\{z:|z|<\varepsilon\}$,
and suppose
\begin{equation}
 R_{\max}/q\le\varepsilon,\qquad t_{\max}/q\le\varepsilon.
 \tag{FTC15}
\end{equation}
The second inequality follows from the first for these original parameters:
$t_{\max}<k/2$ whereas $R_{\max}>2k$.
Pair the $q$ roots of $\chi(c+S)$, with their multiplicities, as $w,-w$.
There are exactly $q/2$ pairs.  Thus
\begin{align}
 |\chi(c+iy)|^2
 &=\prod_{\{w,-w\}}|y^2+w^2|^2
 \le(y^2+R_{\max}^2)^q,\notag\\
 |\chi(c+iqz)|^2
 &\le q^{2q}(2\varepsilon^2)^q
 \qquad (|z|\le\varepsilon).
 \tag{FTC16}
\end{align}
On the segment $1\le z\le2$ the same paired product, with the reverse
triangle inequality, gives
\begin{equation}
 |\chi(c+iqz)|^2
 \ge q^{2q}(1-R_{\max}^2/q^2)^q
 \ge q^{2q}(3/4)^q.
 \tag{FTC17}
\end{equation}
For the numerator factor, the original roots give
\begin{equation}
 |\phi_s(c+iqz)|^2\le q^{2s}(2\varepsilon^2)^s
 \quad (|z|\le\varepsilon),\qquad
 |\phi_s(c+iqz)|^2\ge q^{2s}
 \quad (1\le z\le2).
 \tag{FTC18}
\end{equation}
The factor $(2\varepsilon^2)^s$ is kept in the following stronger bound:
\begin{align}
 \mathcal H_{s,\mathcal I_\varepsilon}
 &\le\kappa_{s,r}^{\mathrm{in}}\mathcal H_{s,\mathcal B_{64}},\notag\\
 \kappa_{s,r}^{\mathrm{in}}
 &\le\frac{2\varepsilon r^2}{\sqrt{\cos1}}
 (2\varepsilon^2)^s\beta^q
 \le\frac{2\varepsilon r^2}{\sqrt{\cos1}}e^{-2q},
 \qquad \beta=\frac43\,2^{-13}e^{2\pi}<e^{-2}.
 \tag{FTC19}
\end{align}
To prove the first inequality, bound $\sigma(qz)$ above by
$\Gamma(1/4)^2/(2\pi\sqrt{\cos1})$ on the inner interval, using (FTC5)
and $e^{-|qz|}\le1$.  Equations (FTC10) and (FTC16)--(FTC18) bound its
integral by
\[
 2\varepsilon q^{2q+2s+1}
 \frac{\Gamma(1/4)^2}{2\pi\sqrt{\cos1}}
 r^2(6+4\varepsilon)^{2(r-1)}
 (2\varepsilon^2)^{q+s}I_Q.
\]
The bulk denominator from $[1,2]$ is at least
$q^{2q+2s+1}c_L(3/4)^qe^{-2\pi q}I_Q$.
The ratio of Gamma constants is again exactly $(\cos1)^{-1/2}$.
Since $6+4\varepsilon<8$ and $r-1\le q$, the polynomial factor is at
most $8^{2q}=2^{6q}$, and
$2\varepsilon^2\cdot2^6=2^{-13}$.  This proves the displayed value of
$\beta$ and the first bound in (FTC19).

The elementary strict estimate on $\beta$ can be certified without decimal
approximations.  One has $\log(4/3)<1/3$ from
$\log(1+x)<x$, and
$\log2=2\sum_{n\ge0}3^{-2n-1}/(2n+1)>2/3$, obtained by integrating the
geometric series for $(1-t^2)^{-1}$ from zero to $1/3$.
Polynomial division gives
\[
 \frac{x^4(1-x)^4}{1+x^2}
 =x^6-4x^5+5x^4-4x^2+4-\frac4{1+x^2}.
\]
Integrating on $[0,1]$ proves
$22/7-\pi=\int_0^1x^4(1-x)^4/(1+x^2)dx>0$.
It follows that
\[
 \log\beta<\frac13-\frac{26}3+\frac{44}7
 =-\frac{43}{21}<-2.
\]
Finally $(2\varepsilon^2)^s\le1$, proving the last bound in (FTC19).

\subsection{Exact determinant transfer after both intervals are included}

Define the complete central set
\begin{equation}
 \mathcal C=\{z:\varepsilon\le|z|\le64\}.
 \tag{FTC20}
\end{equation}
It includes the interval $\varepsilon<|z|<1$.  Except for endpoints of
Lebesgue measure zero, the sets $\mathcal I_\varepsilon$, $\mathcal C$,
and $\mathcal F_{64}$ partition the entire original real integration
line after $y=qz$.  All four matrices in
\begin{equation}
 H_{s,\mathbb R}=H_{s,\mathcal C}
 +H_{s,\mathcal I_\varepsilon}+H_{s,\mathcal F_{64}},\qquad
 0\le H_{s,\mathbb R}-H_{s,\mathcal C}
 \le\kappa_{s,r}H_{s,\mathcal C},\quad
 \kappa_{s,r}=\kappa_{s,r}^{\mathrm{in}}
 +\kappa_{s,r}^{\mathrm{far}}(64,1)
 \tag{FTC21}
\end{equation}
are in the original basis.  The last inequality follows from (FTC12) and
(FTC19), because $\mathcal B_{64}\subset\mathcal C$ and integrals of
nonnegative functions preserve that inclusion.  Thus neither the central
annulus nor any cross term of a polynomial Gram matrix has been discarded.

The matrix $H_{s,\mathcal C}$ is positive definite: $\sigma$ is positive,
the fixed polynomials $\chi\phi_s$ have only finitely many real-line
zeros, and a nonzero polynomial $P$ cannot vanish on an interval.
The Gamma upper bound proves convergence of every entry on the full line.
Let $L$ be the Cholesky factor defined by
$H_{s,\mathcal C}=LL^*$.  The explicit matrix
\[
 A=L^{-1}(H_{s,\mathbb R}-H_{s,\mathcal C})L^{-*}
\]
has $0\le A\le\kappa_{s,r}I$ by (FTC21), and
$H_{s,\mathbb R}=L(I+A)L^*$.  Hence, without changing the original
determinants,
\begin{equation}
 0\le\delta_{s,r}:=
 \log\frac{\det H_{s,\mathbb R}}{\det H_{s,\mathcal C}}
 =\log\det(I+A)\le r\log(1+\kappa_{s,r}).
 \tag{FTC22}
\end{equation}
This argument also proves the same far-only determinant bound if the
inner interval has already been added: its base matrix then dominates
$H_{s,\mathcal B_{64}}$.

At the original high endpoint $N=q+r-1$, define, in precisely the same
relation basis $(\chi,\chi S,\ldots,\chi S^{r-1})$,
\[
 \mathcal R_{N}^{\mathrm{all}}=
 \frac{\det H_{l,\mathbb R}}{\det H_{0,\mathbb R}},\qquad
 \mathcal R_N^{\mathcal C}=
 \frac{\det H_{l,\mathcal C}}{\det H_{0,\mathcal C}}.
\]
Their exact signed relation is
\begin{equation}
 \log\mathcal R_N^{\mathrm{all}}-\log\mathcal R_N^{\mathcal C}
 =\delta_{l,r}-\delta_{0,r}
 \in[-r\log(1+\kappa_{0,r}),\ r\log(1+\kappa_{l,r})].
 \tag{FTC23}
\end{equation}
If the numerator is written in the original Gamma-convolution density
\[
\begin{gathered}
 r_k^0(y)=\frac{(\sqrt{2\pi})^k}{c_{k/4}}
 \frac{|\Gamma(k/4+iy/2)|^2}{2\pi}
 =C_k|f_l(c+iy)|^2\sigma(y),\\
 c_{k/4}=2^{1-k/2}\Gamma(k/2),\quad
 C_k=\frac{(\sqrt{2\pi})^k}{c_{k/4}}4^{-l},
\end{gathered}
\]
the displayed density identity follows by $l$ applications of
$\Gamma(z+1)=z\Gamma(z)$ at $z=1/4+iy/2$.
The recurrence itself follows by integration by parts in Euler's integral,
with both boundary terms zero for $\operatorname{Re}z>0$; squaring the
$l$ factors gives exactly $4^{-l}\prod_j(y^2+t_j^2)$.
Its Gram matrix is exactly $C_kH_{l,A}$,
so its determinant is exactly $C_k^r\det H_{l,A}$.  The term
$r\log C_k$ is present both before and after the restriction and cancels
in (FTC23); its original value is never set to one.

Under (FTC15), both tails have the common bound
\begin{equation}
 \kappa_{s,r}\le\mathcal K_q:=
 \frac1{\sqrt{\cos1}}
 \left(8\varepsilon q^2e^{-2q}+\frac{512q}{59}e^{-19q}\right).
 \tag{FTC24}
\end{equation}
Indeed (FTC15) implies the domain of (FTC14), and $r\le2q$ in (FTC19).
In particular the sum of the two original high-endpoint changes satisfies
\begin{equation}
 \left|\sum_{N\in\{2q-1,2q\}}
 \left(\log\mathcal R_N^{\mathrm{all}}-
 \log\mathcal R_N^{\mathcal C}\right)\right|
 \le(2q+1)\log(1+\mathcal K_q)
 \le(2q+1)\mathcal K_q=o(q).
 \tag{FTC25}
\end{equation}
The two dimensions in this equation remain $q$ and $q+1$, respectively.
The last assertion follows directly from $q^j e^{-2q}\to0$ for each
fixed nonnegative integer $j$, as follows, for example, by retaining
the $(j+1)$-st positive term in the exponential power series.

Finally, the sufficient finite threshold for all of (FTC15)--(FTC25) is
\begin{equation}
 q\ge1024k\sqrt{d+g_\gamma}.
 \tag{FTC26}
\end{equation}
For fixed original parameters this follows from
$k\ge1024\sqrt{d+g_\gamma}$, since $q\ge k^2$.
If a companion central-density Taylor expansion requires the stronger
condition $\tau_\varepsilon\le1/4$, the explicit sufficient threshold is
$q\ge2048k\sqrt{d+g_\gamma}$ instead.  This latter threshold implies
(FTC26); it is not tacitly inferred from equality in (FTC26), which only
gives $\tau_\varepsilon\le1$.  All estimates above retain the original
$\delta,\gamma,m,k,q,r$ and the original Gamma density throughout.
